% Generated draft: requires mathematical and source review.
\documentclass[11pt]{article}
\usepackage[margin=25mm]{geometry}
\usepackage{fontspec}
\setmainfont{Latin Modern Roman}
\usepackage{amsmath,amssymb,mathtools,unicode-math}
\setmathfont{Latin Modern Math}
\usepackage{enumitem,needspace,xurl,hyperref}
\setlength{\parindent}{0pt}
\setlength{\parskip}{5pt}
\setlength{\emergencystretch}{3em}
\newcommand{\N}{\mathbb{N}}
\newcommand{\Z}{\mathbb{Z}}
\newcommand{\Q}{\mathbb{Q}}
\newcommand{\R}{\mathbb{R}}
\newcommand{\C}{\mathbb{C}}
\newcommand{\F}{\mathbb{F}}
\newcommand{\A}{\mathbb{A}}
\newcommand{\PP}{\mathbb{P}}
\newcommand{\E}{\mathbb{E}}
\newcommand{\eps}{\varepsilon}
\newcommand{\dd}{\,\mathrm{d}}
\begin{document}
\section*{UnsolvedMath ID: 30003763}
\textbf{Source problem code:} catalogue entry
\subsection{Definitions and mathematical statement}
The supplied record poses a research programme rather than a single fully specified conjecture. Its base field, local fields, cohomological obstructions, and meaning of the approximation defect are not fixed. The following gives a conditional formulation of its questions without selecting these missing data.

Let \(k\) be a field and, once specified, let \(\Omega\) be a family of places with topological completions \(k_v\). For a \(k\)-variety \(X\), give each \(X(k_v)\) its local topology and put
\[
P_\Omega(X)=\prod_{v\in\Omega}X(k_v),\qquad
\Delta_X:X(k)\longrightarrow P_\Omega(X).
\]
A \(k\)-point means an element of \(X(k)\). Weak approximation with respect to \(\Omega\) means that \(\Delta_X(X(k))\) is dense in \(P_\Omega(X)\) with its product topology. Equivalently, for every finite \(S\subseteq\Omega\), every tuple \((x_v)_{v\in S}\), and every collection of open neighbourhoods \(U_v\subseteq X(k_v)\) of \(x_v\), there exists \(x\in X(k)\) whose image lies in \(U_v\) for every \(v\in S\).

Suppose a specified collection of cohomological obstructions defines an admissible locus \(\mathcal C_\Omega(X)\subseteq P_\Omega(X)\) containing the diagonal rational points. In this framework, the rational-point question asks for classes of varieties for which
\[
\mathcal C_\Omega(X)\ne\varnothing\quad\Longrightarrow\quad X(k)\ne\varnothing.
\]
The corresponding approximation question asks for classes for which every admissible tuple can be approximated by rational points, that is,
\[
\mathcal C_\Omega(X)\subseteq
\overline{\Delta_X(X(k))}.
\]
These are formulation templates: the source does not identify \(\mathcal C_\Omega(X)\), or say whether the obstructions are instead defined on an adelic or another restricted space of local points. No particular cohomological obstruction is selected here.

The source also asks how such cohomological methods relate to patching over function fields. It does not specify the function fields, models, patch fields, or the comparison to be established.

Finally, it asks for the defect of weak approximation for arbitrary noncommutative algebraic groups. For a \(k\)-algebraic group \(G\) and finite \(S\subseteq\Omega\), one possible measure of failure is the space of right cosets
\[
P_S(G)=\prod_{v\in S}G(k_v),\qquad
H_S(G)=\overline{\Delta_{G,S}(G(k))},\qquad
D_S(G)=P_S(G)/H_S(G).
\]
Here \(\Delta_{G,S}\) is the diagonal map. The closure \(H_S(G)\) is a subgroup; its coset space is a quotient group only if it is normal. Weak approximation holds precisely when \(H_S(G)=P_S(G)\) for every finite \(S\). The source does not specify whether these coset spaces, a global analogue, or another invariant is its intended defect. Thus this construction illustrates a possible precise interpretation, not an identification asserted by the source. The intended classes of groups and any smoothness, connectedness, affineness, or reductivity assumptions remain unspecified.
\subsection{Short English statement}
The problem asks when cohomological obstructions completely explain failure of rational points and weak approximation, how these methods compare with patching over function fields, and how to describe failure of weak approximation for noncommutative algebraic groups. The field, local data, obstruction theory, and intended defect must be specified before this becomes a definite conjecture.
\subsection{Sources}
\begin{itemize}
\item Supplied local UnsolvedMath registry record, permanent ID \(30003763\), source problem code \texttt{catalogue entry}, title: ``Cohomological Obstructions to Rational Points and Weak Approximation.'' Collection: Oberwolfach Reports Open Problems. No individual proposer is supplied; the proposed year is recorded as 2018.
\item Supplied report attribution: \emph{Mini-Workshop: Arithmetic Geometry and Symmetries around Galois and Fundamental Groups (2019)}, Oberwolfach Reports, pp.\ 1021--1080, DOI \texttt{10.4171/owr/2018/17}. \url{https://doi.org/10.4171/owr/2018/17}. The literature-source list labels this reference ``OWR 2018 report.'' The supplied background describes these reports as records of mathematical workshops and their open-problem discussions.
\item Supplied literature reference: ``Harari--Szamuely, local-global principles.'' \url{https://doi.org/10.1090/jag/661}. No fuller bibliographic details or precise supporting theorem are supplied.
\item Literature assessment dated exactly \texttt{2026-08-21}: the supplied curation labels the problem open and states that patching and cohomological duality solve important torus and selected reductive cases over function fields, while the requested general classification of local-global obstructions and weak approximation for arbitrary noncommutative groups remains open. This is supplied dated curation, not an independently verified literature assessment. The background expressly requests independent verification of time-sensitive status and cited claims.
\end{itemize}
\end{document}
